\chapter{Mapa de cobertura y estrategia}
\section{Qué establece el documento oficial}
El PDF aportado describe 140 reactivos disciplinares y 60 transversales dentro de una estructura general de 200; además menciona reactivos piloto. No debe confundirse esa composición con el tamaño del banco de práctica de esta guía. El documento señala tres opciones por pregunta y presenta cuestionamiento directo, completamiento, ordenamiento y relación de elementos, además de estímulos compartidos \cite[págs. 8--11 y 35--36]{ceneval}.

En la sección disciplinar, los pesos publicados son: hardware 49, redes 41 y software 50. Aquí se ofrecen cien preguntas por área, tal como se solicitó. Por ello, el banco refuerza cada área por igual y no reproduce la ponderación global del examen. Dentro de cada una, la distribución aproxima proporcionalmente sus subáreas.

\section{Matriz de trazabilidad}
Los desempeños siguientes son síntesis editoriales de los descriptores, no transcripciones ni una lista oficial exhaustiva de temas. En cada sección disciplinar hay al menos tres ejemplos desarrollados, numerados E1, E2 y E3. El intervalo de preguntas permite localizar práctica de la misma subárea.
\small
\begin{longtable}{>{\raggedright\arraybackslash}p{0.17\textwidth}>{\raggedright\arraybackslash}p{0.35\textwidth}>{\raggedright\arraybackslash}p{0.17\textwidth}>{\raggedright\arraybackslash}p{0.18\textwidth}}
\toprule Subárea & Desempeño que se prepara & Teoría/ejemplos & Reactivos\\\midrule\endhead
1.1 Circuitos & Analizar e integrar circuitos analógicos y digitales & \ref{sec:h1}, E1--E4 & H-001--H-022\\
1.2 Arquitectura & Interpretar números, memoria, organización e instrucciones & \ref{sec:h2}, E1--E4 & H-023--H-044\\
1.3 Embebidos & Implementar firmware, interfaces, comunicaciones y periféricos & \ref{sec:h3}, E1--E4 & H-045--H-071\\
1.4 Automatización & Modelar señales y sistemas, seleccionar sensores y aplicar control & \ref{sec:h4}, E1--E4 & H-072--H-100\\
2.1 Diseño & Seleccionar topologías, medios, componentes y direccionamiento & \ref{sec:r1}, E1--E4 & R-001--R-027\\
2.2 Administración & Configurar, mantener, monitorear y diagnosticar servicios y red & \ref{sec:r2}, E1--E4 & R-028--R-059\\
2.3 Seguridad & Identificar riesgos y evaluar controles, criptografía y funcionamiento & \ref{sec:r3}, E1--E4 & R-060--R-100\\
3.1 Ingeniería & Analizar requisitos, diseñar, probar, mantener y estimar productos & \ref{sec:s1}, E1--E4 & S-001--S-032\\
3.2 Software base & Interpretar programas, compiladores y gestión del sistema operativo & \ref{sec:s2}, E1--E4 & S-033--S-058\\
3.3 Datos y algoritmos & Usar estructuras, recursividad, transacciones e IA básica & \ref{sec:s3}, E1--E5 & S-059--S-100\\
Lectura: estudio & Identificar, interpretar y evaluar reseña e informe & \ref{sec:l1} & L-001--L-012\\
Lectura: literario & Interpretar cuento y ensayo sin exceder su evidencia & \ref{sec:l2} & L-013--L-024\\
Lectura: social & Aplicar criterios de convocatoria y evaluar nota & \ref{sec:l3} & L-025--L-030\\
Redacción: estudio & Seleccionar textos adecuados, coherentes y normativos & \ref{sec:e1} & E-001--E-015\\
Redacción: social & Adecuar avisos, solicitudes y notas al propósito & \ref{sec:e2} & E-016--E-030\\
\bottomrule
\end{longtable}\normalsize
\fuentes{ceneval}
\section{Método de resolución y control del avance}
Antes de calcular, identifique qué se pregunta, el modelo y las unidades. Separe datos de supuestos: una caché idealizada, un ADC por intervalos o una red convencional pueden tener reglas distintas de un equipo específico. Dibuje el circuito, el flujo o la tabla mínima que haga visible el problema.

Después de obtener un resultado, compruebe escala y límites: una resistencia paralela positiva debe ser menor que la menor rama; una utilización no puede superar el cien por ciento si se define como fracción de capacidad; una probabilidad debe estar entre cero y uno. En código, trace variables y pruebe extremos. En lectura, localice evidencia antes de elegir.

Si dos opciones parecen correctas, revise los calificadores: \emph{siempre}, \emph{necesariamente}, \emph{solo}, \emph{mínimo} y \emph{bajo estos supuestos}. Una afirmación generalmente plausible puede ser incorrecta para el caso definido. Al finalizar, explique por qué descartó cada alternativa; esa justificación revela lagunas que un acierto por azar oculta.

Use una bitácora con identificador, respuesta, certeza, tiempo, causa del error y fecha de revisión. Reparta el repaso en sesiones espaciadas; alterne ejercicios de áreas diferentes una vez dominados los procedimientos básicos. Reserve una sesión final para preguntas sin consultar apuntes y revise después los errores con el solucionario.
